% appendix_bayesopt.tex — methodology appendix for the upcoming paper.
% Standalone \section; drop into the main .tex or \input{} it.
% Citations use authoryear keys; add these to the .bib if absent:
%   vdbosch2019dynamite (DYNAMITE method paper), zhu2018crcut,
%   eriksson2019turbo, elgammal2023gpry, li2026sale, leclercq2018bolfi.
%
% Numbers below come from the runs documented in
% dev_notes/bayesopt_generator.md (June 2026 NGC6278 comparison; August
% 2026 surface benchmarks and tuning study). Update them if the final
% production numbers change.

\section{Bayesian-optimization parameter search}
\label{app:bayesopt}

\subsection{Problem structure}
Triaxial Schwarzschild models separate into a cheap inner problem and an
expensive outer one \citep{vdbosch2019dynamite}. Given a trial potential,
the orbit-library weights that best reproduce the observed kinematics
follow from a bounded non-negative least-squares solve, which takes
seconds. The outer problem, finding the potential parameters that
minimize the kinematic $\chi^2$, requires one orbit-library integration
per trial potential and takes hours on a multicore node. Classical
implementations search the outer parameters with a grid walk
\citep{zhu2018crcut}: each step proposes axis-aligned neighbours of the
current best point, so the search cost grows with the number of free
parameters and the walk can stall when all neighbours are already
evaluated.

The generator described here, \texttt{BayesOptGenerator}, replaces the
walk with a Gaussian-process surrogate of $\chi^2$ over the normalized
parameter cube. The surrogate is refit after every batch of models, and
the next batch is chosen where the surrogate predicts improvement. The
approach follows established surrogate-based inference practice
\citep{leclercq2018bolfi,elgammal2023gpry,li2026sale}; the published
Schwarzschild-modelling literature searches the outer parameters with
grids or walks.

\subsection{Algorithm}
\label{app:bayesopt:algo}
Each iteration proceeds in four steps.

\paragraph{(i) Training data.} All evaluated models with finite
$\chi^2$ are collected from the run table. Parameters are normalized to
the unit cube using their prior bounds; logarithmic parameters
($M_\mathrm{BH}$, halo mass) enter in log space. Rows outside the
current bounds, which occur when bounds are tightened between runs, are
clipped and counted in a warning.

\paragraph{(ii) Surrogate.} A BoTorch \texttt{SingleTaskGP} with
outcome standardization and an anisotropic Mat\'ern~5/2 kernel is fit by
empirical Bayes; above 300 training rows the code switches to a
variational GP for scaling. The GP models $-\chi^2$.

\paragraph{(iii) Batch selection.} The next $q$ candidates maximize
$\log q$-Expected Improvement \citep[qLogEI;][]{eriksson2019turbo}
jointly over the cube. Three modifications matter in practice.
First, exploration is annealed: the indicator temperature $\eta$ of
qLogEI decays linearly from $0.1$ to $10^{-3}$ over the first ten
batches, so early batches favour uncertain regions and later batches
exploit. Second, a fraction $\lceil q/4\rceil$ of each batch is drawn
not from the acquisition optimizer but by rejection sampling from
$\exp(\mu_\mathrm{GP}/\tau)$ over the feasible set, with temperature
$\tau$ annealed from $1.0$ by a factor $0.7$ per batch
\citep[the annealed objective of][]{li2026sale}. These members keep the
batch diverse where joint optimization would homogenize it. Third,
after selection, candidates are deduplicated on the grid cells of the
non-mass parameters: two candidates in the same cell would integrate
the same orbit library, and the freed slots are refilled with
quasi-random draws. Pairs of candidates that share a cell but differ in
$M_\mathrm{L}$ are kept, since the orbit library is then reused at no
integration cost.

\paragraph{(iv) Feasibility.} Triaxiality requires $p\ge q$,
$u \ge \max(q/q_\mathrm{obs}, p)$ and $u \le \min(p/q_\mathrm{obs},1)$
for the intrinsic flattenings. When all three are free, these enter the
acquisition optimization as differentiable inequality constraints.
When only a subset is free, which is the common case (the NGC~5139 run
of Section~\ref{sec:production} fixes $u$), the fixed values are
substituted into the constraints and the free axes are projected onto
the feasible set before any evaluation is proposed.

\paragraph{Warm-up and warm start.} The first $n_\mathrm{init}$
candidates come from a scrambled Sobol sequence, or, optionally, from an
axial design around a user-supplied point. When the generator is
started in an output directory that already contains evaluated models,
for example after a grid-walk search, the Sobol phase is skipped and
the GP trains on the full history. The axial design then centres on the
best historical model. This makes switching from a grid walk to
Bayesian optimization mid-project free of redundant evaluations.

\paragraph{Convergence diagnostics.} Alongside the standard
plateau criteria, the generator tracks how often the GP predicts a new
model's $\chi^2$ within a relative tolerance of the realized value
\citep[the CorrectCounter of][]{elgammal2023gpry}. A run of
$\max(4,\lceil d/2\rceil)$ accurate predictions indicates that the
surrogate has resolved the landscape; we use this as a diagnostic and
not as a stopping rule. A single trust region
\citep[TuRBO-style;][]{eriksson2019turbo} can additionally be enabled,
restricting acquisition to a shrinking box around the incumbent once
the mean distance to the $k$ nearest evaluated points falls below a
threshold fraction of the box diagonal.

\subsection{Validation protocol}
\label{app:bayesopt:valid}
We validate on three levels. (a) Unit tests cover the feasibility
projection for all free/non-free subsets of $(q,p,u)$, warm-start
behaviour, batch deduplication, and each acquisition option; a fuzz
suite asserts the feasibility invariants over randomized parameter
spaces. (b) Synthetic $\chi^2$ surfaces calibrated to the ranges of our
galaxy sample encode the structures that make these searches hard: a
mass--flattening valley, a black-hole bowl over a correlated shape
band, and a curved constant-mass degeneracy valley with a secondary
basin. On the coupled four-parameter surface, Bayesian optimization
reaches the convergence threshold in a median of 42 fresh models
against 116 for the grid walk over six seeds, with a better final
$\chi^2$ (6013 versus 6070); on single-valley surfaces the grid walk
remains competitive, which bounds the recommendation. A hard acceptance
test requires the generator to reach the threshold in fewer models than
the grid walk and keeps surrogate overhead below ten seconds per model.
(c) On NGC~6278 with three free parameters and equal model budgets, the
generator reaches $\chi^2=5814$ where the grid walk reaches 6352, and
finds its optimum between grid points.

\subsection{Recommended settings}
For the production runs we use batch size $q=4$, which the tuning study
identifies as the setting with the largest effect at $d=4$ (median
models-to-threshold 26 versus 41 for $q=8$), consistent with the
batch-size guideline of \citet{elgammal2023gpry}; $n_\mathrm{init}=8$
Sobol points; the annealed $\eta$ schedule; $q/4$ tempered-posterior
members; and trust regions enabled. We advise running a grid walk first
and switching generators in the same output directory, so that the GP
starts from the walk's full evaluation history.
